\scanpage{17}

et une \deleted{\unsure{ho.|la}} forme quadratique de définition de la structure \deleted{quadratique} \inserted{de} similitude

\[
Q_\Sigma:E\longrightarrow\Delta^\vee\overset{\mathrm{déf}}{=}\Delta' \;(\simeq\mathcal A/\mathcal O\otimes L^\vee)
\]

\deleted{Ceci posé, soit} Plus généralement, soit $\Delta'$ un Module inversible,

\[
Q:E\longrightarrow\Delta'\qquad(\text{i.e. }\Delta^{\prime\vee}\longrightarrow\underline{\mathrm{Quad}}(E)\quad\text{un hom. de Modules})
\]

une forme quadratique, d'où une structure de similitudes $\Sigma=\Sigma_Q$ sur $E$, \unsure{correspond.|correspond\textsuperscript{t}|correspondant} à un sous-anneau $\mathcal A\unsure{=|\supset}\mathcal A_\Sigma$ de $\underline{\mathrm{End}}(E)$.

Soit $x$ une section de $E$. Conditions équivalentes

a) $\{x\}$ est un générateur de $E$ comme $\mathcal A$-Module, i.e. \inserted{$\lambda\mapsto\lambda\cdot x$,} $\mathcal A\longrightarrow\deleted{\unsure{\mathcal A}}E$, est épi

a$'$) \deleted{$Q(x)$ \unsure{es|est}} $\{x\}$ est une base de $E$ comme $\mathcal A$ Module, i.e. \inserted{$u\mapsto u\cdot x$:} $\mathcal A\longrightarrow E$ est un iso

b) $\{Q(x)\}$ est une base de $\deleted{\unsure{E|L}}\Delta'$

b$'$) (Si $\Delta'=\underline{\mathcal O}$) $Q(\unsure{u|x})\in\Gamma\underline{\mathcal O}_X$ est \uline{inversible}.

c) $\forall s\in X$, $Q(s)\ne0$

\deleted{\uline{Cor. 2} Soit $Q$}

d) (Si $u$ est un gén. loc. de l'Alg. $\mathcal A$) Pour \deleted{tout} $s\in X$, $x(s)\in E(s)$ est \uline{non vect.} \uline{propre} de $u(s)\in\mathrm{End}\,E(s)$
% archive sheet number: 16
